\documentclass[10pt,a4paper]{amsart}
\usepackage[margin=19mm]{geometry}
\usepackage{amsmath,amssymb,mathtools}
\usepackage[scr=rsfs,cal=euler]{mathalfa}
\usepackage{xfrac}
\usepackage[unicode,hidelinks,bookmarks=false]{hyperref}
\newcommand{\ie}{i.e.\ }
\newcommand{\cf}{cf.\ }
\newcommand{\resp}{respectively\ }
\newcommand{\Subsection}{\subsection}
\providecommand{\nobreakdash}{\nobreak\hbox{-}}
\setlength{\emergencystretch}{2em}
\multlinegap0pt
\usepackage{amssymb}
\usepackage{xcolor}
\usepackage{tcolorbox}
\usepackage[shortlabels]{enumitem}
\usepackage{float}
\newtheorem{proposition}{Proposition}
\title{Kernel-Gradient Drifting Models}
\author{Maria Esteban-Casadevall, Jorge Carrasco-Pollo, Max Welling, Jan-Willem van de Meent, Erik J. Bekkers, Floor Eijkelboom}
\date{Source: arXiv:2605.10727v1 (CC BY 4.0)}
\begin{document}
\maketitle

\noindent\emph{Source and licence.} Kernel-Gradient Drifting Models. Version arXiv:2605.10727v1 (CC BY 4.0), distributed by arXiv under the Creative Commons Attribution 4.0 International licence, http://creativecommons.org/licenses/by/4.0/, as recorded in the arXiv licence metadata for this exact version and shown on the abstract page https://arxiv.org/abs/2605.10727v1. Full licence notice: \texttt{licenses/arxiv-2605.10727-CC-BY-4.0.txt}.\par\smallskip

\noindent\textit{Editorial scope.} Counted unit: the article's proposition prop:score\_ratio\_matching (source0.tex lines 434--458). The article's complete original proof of it is delivered with it (source0.tex lines 1312--1326, one \texttt{proof} environment of 15 lines, which the article places away from the statement). No mathematics is written, repaired or re-derived: the statement and the proof are the article's own lines, in the article's own notation, and no retained line is substituted or re-flowed.  Retained uncounted as the source-local context the proof needs: lines 357--386.  Nothing was omitted from any retained span, so the package carries no omission record.  No retained line contains a citation, so the wrapper carries no bibliography and no bibliography entry.  No retained line contains a figure environment, an image-inclusion command or drawing code, so no image is delivered and the package declares no asset dependency.  Editorial formatting only, in addition: an AMS \texttt{amsart} preamble that restates the article's own declarations and macros used by the retained lines, namely the environments \texttt{proposition} on separate counters, together with the standard packages those lines need.  The retained environments print in this document as Proposition 1 in order of appearance, whereas the article prints the same items with the numbering of its own full document.  The complete native source of the article as distributed in the arXiv e-print for this exact version is bundled unchanged as \texttt{sources/200314/source0.tex}.
\begin{equation}
\label{eq:gradient_drifting}
V^\nabla_{p,q_\theta}(x)
=
G_{p}^+(x) - G_{q_\theta}^-(x)
=
\frac{
\mathbb{E}_{y \sim p}
\!\left[
\nabla_x k(x,y)
\right]
}{
\mathbb{E}_{y \sim p}
\!\left[
k(x,y)
\right]
}
-
\frac{
\mathbb{E}_{x' \sim q_\theta}
\!\left[
\nabla_x k(x,x')
\right]
}{
\mathbb{E}_{x' \sim q_\theta}
\!\left[
k(x,x')
\right]
}.
\end{equation}

\begin{proposition}[Kernel-gradient drifting is smoothed score-ratio matching]
\label{prop:score_ratio_matching}
Let \(k\) be a positive, normalized kernel differentiable in its first argument, and assume that
differentiation may be exchanged with integration. Then the kernel-gradient drift
defined in \eqref{eq:gradient_drifting} satisfies
\[
V^\nabla_{p,q_\theta}(x)
=
\nabla_x \log
\frac{\widehat{p}_k(x)}
     {\widehat{q}_{\theta,k}(x)}.
\]
Consequently,
\[
\mathcal{L}^{\nabla}_\mathrm{drift}(\theta)
=
\eta^2
\mathbb{E}_{x\sim q_\theta}
\left\|
\nabla_x \log \widehat{p}_k(x)
-
\nabla_x \log \widehat{q}_{\theta,k}(x)
\right\|_2^2.
\]
\end{proposition}

\begin{proof}
By Proposition~\ref{prop:score_ratio_matching}, \(V^\nabla_{p,q}=0\) implies
\[
\nabla_x \log \widehat{p}_k(x)
=
\nabla_x \log \widehat{q}_k(x).
\]
As $k$ is a normalized kernel, this implies
\[
\widehat{p}_k(x)
=
\widehat{q}_k(x).
\]
Since \(k\) is characteristic, this implies \(p=q\).
\end{proof}

\end{document}
